% !TeX root = ../main.tex

% ==================================================
% 1.3. sys.path, pathlib
% ==================================================

\subsection{\zsh{sys.path, pathlib}}

Python에서 "경로"라는 값이 보인다고 해서 모두 같은 목적의 경로는 아니다.

무엇을 찾기 위한 경로인지 구분해야 한다.

\subsubsection{먼저 실습용 프로젝트를 하나로 고정한다}

앞으로 이 구조를 계속 사용한다.

\begin{dhkdir}[in=9]
my_ml_project/
>data/
>>train.csv
>>test.csv
>model/
>notebooks/
>>experiment.ipynb
>src/
>>train.py
>>uitls.py
>README.md
\end{dhkdir}



Mac에서 실제 프로젝트 위치가 다음이라고 가정하자.

\begin{dhkzsh}
/Users/user/Desktop/my_ml_project
\end{dhkzsh}

따라서 \zsh{train.py}의 절대 경로는

\begin{dhkzsh}
/Users/user/Desktop/my_ml_project/src/train.py
\end{dhkzsh}

이다.

이제 세 가지를 보자

\begin{dhksmry}[centered]
Path.cwd()

__file__

sys.path
\end{dhksmry}

겉으로는 셋 다 \zsh{/Users/...}같은 경로가 튀어나올 수 있다.

하지만 질문 자체가 다르다
\begin{dhknote}
\begin{dhksmry}[centered]
Path.cwd()
지금 어디에서 작업 중인가?


__file__
지금 실행하는 이 파일은 어디 있는가?


sys.path
mport 할 때 Python이 어디를 검색하는가?
\end{dhksmry}
\end{dhknote}

\subsubsection{\zsh{Path.cwd()} - 나는 지금 어디서 작업하고 있는가?}

\begin{dhkcode}
from pathlib import Path

print(Path.cwd())
\end{dhkcode}

\zsh{CWD}는 \zsh{Current Working Directory}, 즉 \textbf{현재 작업 디렉터리}다.

터미널에서:

\begin{dhkzsh}
/Users/user/Desktop/my_ml_project %
\end{dhkzsh}

인 상태에서:

\begin{dhkzsh}
python src/train.py
\end{dhkzsh}

를 실행했다고 하자.

그러면 보통:

\begin{dhkcode}
Path.cwd()
\end{dhkcode}

결과는:

\begin{dhkzsh}
/Users/user/Desktop/my_ml_project
\end{dhkzsh}

이다.

\zsh{train.py}는 여기 있다.

\begin{dhkzsh}
/Users/user/Desktop/my_ml_project/src/train.py
\end{dhkzsh}

그런데 \zsh{CWD}는:

\begin{dhkzsh}
/Users/user/Desktop/my_ml_project
\end{dhkzsh}

이며, 둘은 다르다.

\subsubsection{\zsh{__file__} - 이 Python 파일 자체는 어디 있는가?}

\zsh{train.py}안에 다음 코드를 넣어보자.

\begin{dhkcode}
print(__file__)
\end{dhkcode}

그러면 실행 중인 \textbf{해당 Python 파일의 경로}를 알 수 있다.

더 확실하게 절대 경로 형태로 보고 싶다면:

\begin{dhkcode}
from pathlib import Path

print(Path(__file__).resolve())
\end{dhkcode}

결과는:

\begin{dhkzsh}
/Users/user/Desktop/my_ml_project/src/train.py
\end{dhkzsh}

같은 식이다.

따라서:

\begin{dhkcode}
Path.cwd()

/Users/user/Desktop/my_ml_project
\end{dhkcode}

와:

\begin{dhkcode}
Path(__file__).resolve()

/Users/user/Desktop/my_ml_project/src/train.py
\end{dhkcode}

는 전혀 다른 것을 말하고 있다.

\subsubsection{이 차이를 그림으로 보면 간단하다}

현재:

\begin{dhkdir}
/
>Users/
>>user/
>>>Desktop/
>>>>my_ml_project - CWD
\end{dhkdir}

여기에서:

\begin{dhkzsh}
python src/train.py
\end{dhkzsh}

를 실행했다.

그러면:

\begin{dhkdir}
my_ml_project <- Path.cwd()
>data/
>src/
>>train.py  <- __file__
\end{dhkdir}

즉:

\begin{dhknote}
\begin{dhksmry}[centered]
CWD
내가 프로그램을 실행하고 있는 작업 기준 위치

__file__
실행되는 이 Python 파일 자체의 위치
\end{dhksmry}
\end{dhknote}

이다.

\subsubsection{그래서 \zsh{__file__}을 이용하면 파일 기준 경로를 만들 수 있다}

이 패턴은 실제 프로젝트에서 매우 유용하다

\zsh{train.py}가:

\begin{dhkzsh}
my_ml_project/src/train.py
\end{dhkzsh}

에 있다.

\begin{dhkcode}
from pathlib import Path

current_file = Path(__file__).resolve()

print(current_file)
\end{dhkcode}

결과:

\begin{dhkzsh}
/Users/user/Desktop/my_ml_project/src/train.py
\end{dhkzsh}

그런데 우리는 파일이 아니라 src 디렉터리를 원한다고 해보자.

\begin{dhkcode}
current_dir = current_file.parent
\end{dhkcode}

결과:

\begin{dhkzsh}
/Users/user/Desktop/my_ml_project/src
\end{dhkzsh}

한 번 더:

\begin{dhkcode}
project_root = current_file.parent.parent
\end{dhkcode}

결과:

\begin{dhkzsh}
/Users/user/Desktop/my_ml_project
\end{dhkzsh}

가 된다.

따라서:

\begin{dhkcode}
data_path = project_root / "data" / "train.csv"
\end{dhkcode}

하면:

\begin{dhkzsh}
/Users/user/Desktop/my_ml_project/data/train.csv
\end{dhkzsh}

가 된다.

전체 코드:

\begin{dhkcode}
from pathlib import Path

current_file = Path(__file__).resolve()
project_root = current_file.parent.parent
data_path = project_root / "data" / "train.csv"
print(data_path)
\end{dhkcode}

이 코드는 \zsh{CWD}에 덜 의존한다.

\subsubsection{왜 그런지 살펴보자}

\zsh{train.py}를 다음처럼 작성하자.

\begin{dhkcode}
from pathlib import Path

print("CWD :", Path.cwd())
print("FILE:", Path(__file__).resolve())
\end{dhkcode}

프로젝트:

\begin{dhkdir}
my_ml_project/
>src/
>>train.py
\end{dhkdir}


이제 실행 방법을 바꿔보자.

\subsubsection{실험 A - 프로젝트 \zsh{root}에서 실행}

터미널:

\begin{dhkzsh}
my_ml_project %
\end{dhkzsh}

여기서:

\begin{dhkzsh}
python src/train.py
\end{dhkzsh}

결과:

\begin{dhkzsh}
CWD : /Users/user/Desktop/my_ml_project
FILE: /Users/user/Desktop/my_ml_project/src/train.py
\end{dhkzsh}

즉:

\begin{dhkzsh}
CWD
my_ml_project/

FILE
my_ml_project/src/train.py
\end{dhkzsh}

\subsubsection{실험 B - \zsh{src} 안으로 들어가서 실행}

터미널에서:

\begin{dhkzsh}
cd src
\end{dhkzsh}

그러면:

\begin{dhkzsh}
my_ml_project/src %
\end{dhkzsh}

이제:

\begin{dhkzsh}
python train.py
\end{dhkzsh}

결과:

\begin{dhkzsh}
CWD : /Users/user/Desktop/my_ml_project/src
FILE: /Users/user/Desktop/my_ml_project/src/train.py
\end{dhkzsh}

비교해보자.

\textbf{실험 A}

\begin{dhkzsh}
CWD = my_ml_project
FILE = my_ml_project/src/train.py
\end{dhkzsh}

\textbf{실험 B}

\begin{dhkzsh}
CWD = my_ml_project/src
FILE = my_ml_project/src/train.py
\end{dhkzsh}

\textbf{\zsh{__file__}은 그대로인데 \zsh{CWD}가 바뀌었다}.

이게 핵심이다.

\subsubsection{그래서 상대 경로의 결과가 달라진다}

\zsh{train.py}에 다음 코드가 있다고 하자.

\begin{dhkcode}
from pathlib import path

data = Path("data/train.csv")
print(data.resolve())
\end{dhkcode}

\textbf{프로젝트 root에서 실행하면}

\begin{dhkzsh}
my_ml_project %    
\end{dhkzsh}

에서:

\begin{dhkzsh}
python src/train.py
\end{dhkzsh}

하면:

\begin{dhkzsh}
CWD
/Users/user/Desktop/my_ml_project
\end{dhkzsh}

이므로:

\begin{dhkzsh}
data/train.csv
\end{dhkzsh}

는:

\begin{dhkzsh}
/Users/user/Desktop/my_ml_project/data.train.csv
\end{dhkzsh}

가 된다.

\subsubsection{그런데 \zsh{src}에서 실행하면?}

\begin{dhkzsh}
my_ml_project/src %
\end{dhkzsh}

에서:

\begin{dhkzsh}
python train.py
\end{dhkzsh}

하면 \zsh{CWD}가:

\begin{dhkzsh}
/Users/user/Desktop/my_ml_project/src
\end{dhkzsh}

다.

따라서:

\begin{dhkcode}
Path("data/train.csv")
\end{dhkcode}

는:

\begin{dhkzsh}
/Users/user/Desktop/my_ml_project/src/data/train.csv
\end{dhkzsh}

를 의미한다.

그런데 실제 파일은:

\begin{dhkzsh}
/Users/user/Desktop/my_ml_project/data/train.csv
\end{dhkzsh}

에 있다.

그러면 \zsh{FileNotFoundError}가 발생할 수 있다.

\subsubsection{여기서 아주 중요한 결론}

많은 초보자가 이렇게 생각한다.

\begin{dhkcode}
Path("data/train.csv")
\end{dhkcode}

\textbf{Python 파일이 있는 위치에서 \zsh{data/train.csv}를 찾겠지?}

아니다.

기본적으로 상대 파일 경로를 해석할 때 중요한 기준은:

\begin{dhksmry}[centered]
.py 파일 위치
\end{dhksmry}

가 아니라:

\begin{dhksmry}[centered]
Current Working Directory, CWD
\end{dhksmry}

이다.

그래서

\begin{dhksmry}[centered]
같은 train.py
+
같은 Path("data/train.csv")
\end{dhksmry}

여도 \textbf{어디에서 프로그램을 실행했느냐에 따라 가리키는 곳이 달라질 수 있다}.

\subsubsection{그러면 파일 위치를 기준으로 만들고 싶으면?}

그때 \zsh{__file__}을 이용할 수 있다.

\begin{dhkcode}
from pathlib import Path

project_root = Path(__file__).resolve().parent().parent()
data_path = project_root / "data" / "train.csv"
\end{dhkcode}

구조적으로:

\begin{dhkbranch}[in=8]
node main :: train.py
arrow main
node main :: Path(__file__)
arrow main
node main :: .../my_ml_project/src/train.py
line main :: .parent
arrow main
node main :: .../my_ml_project/src
line main :: .parent
arrow main
node main :: .../my_ml_project
line main :: / "data"
arrow main
node main :: .../my_ml_project/data
line main :: / "train.csv"
arrow main
node main :: .../my_ml_project/data/train.csv
\end{dhkbranch}

이제 실행 위치가 바뀌어도 프로젝트 구조 자체가 유지된다면 같은 파일을 가리킬 수 있다.

\subsubsection{이제 세 번재인 \zsh{sys.path}}

여기서 이름 때문에 혼란이 생긴다.

\begin{dhkcode}
import sys

print(sys.path)
\end{dhkcode}

결과는 대략:

\begin{dhkzsh}
[
'/Users/user/Desktop/my_ml_project/src',
'/Users/user/anaconda3/envs/transformer/lib/python310.zip',
'/Users/user/anaconda3/envs/transformer/lib/python3.10',
'/Users/user/anaconda3/envs/transformer/lib/python3.10/site-packages',
...
]
\end{dhkzsh}

이건 \zsh{train.csv}를 찾기 위한 목록이 아니다.

Python이:

\begin{dhkcode}
import numpy
import pandas
import transformers
\end{dhkcode}

같은 \textbf{모듈/패키지를 import할 때 검색하는 위치}와 관련된 목록이다.

즉:
\begin{dhknote}
\begin{dhksmry}[centered]
Path.cwd()
작업 위치

__file__
이 Python 파일 위치

sys.path
import 검색 위치들
\end{dhksmry}
\end{dhknote}
이다.

\subsubsection{\zsh{sys.path}는 하나의 경로가 아니라 목록이다}

이 차이도 중요하다.

\begin{dhkcode}
Path.cwd()
\end{dhkcode}

는 하나의 현재 작업 위치를 나타낸다.

반면:

\begin{dhkcode}
sys.path
\end{dhkcode}

는 여러 검색 위치를 담고 있다.

개념적으로:

\begin{dhkdir}
sys.path
>위치 A
>위치 B
>위치 C
>site-packages
>...
\end{dhkdir}

Python이:

\begin{dhkcode}
import pandas
\end{dhkcode}

를 만나면 이 검색 체계를 이용해 pandas를 찾는다.

그래서 Conda 환경의:

\begin{dhkzsh}
.../envs/transformer/lib/python3.10/site-packages
\end{dhkzsh}

같은 위치가 중요해진다.

그곳에 transformer, pandas, numpy 등이 설치되어 있기 때문이다.

\subsubsection{그래서 \zsh{sys.path}와 파일 읽기는 분리해야 한다}

예를 들어:

\begin{dhkcode}
import pandas as pd

df = pd.read_csv("data/train.csv")
\end{dhkcode}

에서 \zsh{data/train.csv}를 찾는 문제와 pandas를 찾는 문제는 다르다.

전자는 파일 시스템 + 파일 경로 문제다.

후자는 Python import system + sys.path 문제다.

\subsubsection{지금까지를 하나의 그림으로 통합}

\begin{dhkdir}
Python 실행
> Path.cwd()
>> 작업 기준점 -> 상대 파일 경로 해석 등에 중요
> __file__
>> 파일 위치 -> 프로젝트 root 계산 등에 사용
> sys.path
>> import 검색 -> numpy, pandas, transformer
\end{dhkdir}

\subsubsection{\zsh{pathlib.Path}를 제대로 다뤄보자}

여기서부터는 실제 도구 사용이다.

우선:

\begin{dhkcode}
from pathlib import Path
\end{dhkcode}

한다.

\subsubsection{\zsh{Path}객체 만들기}

\begin{dhkcode}
p = Path("data/train.csv")
\end{dhkcode}

현재로서는 파일을 읽은 게 아니다.

그냥:
\begin{dhksmry}[centered]
data/train.csv라는 경로를 나타내는 객체
\end{dhksmry}

를 만든 것이다.

이 차이가 중요하다.

\begin{dhkcode}
Path("data/train.csv")
\end{dhkcode}

\textbf{한다고 파일이 열리는 것도 아니고 생성되는 것도 아니다}.

\subsubsection{현재 작업 디렉터리}

\begin{dhkcode}
Path.cwd()
\end{dhkcode}

예:

\begin{dhkzsh}
/Users/user/Desktop/my_ml_project
\end{dhkzsh}

\subsubsection{홈 디렉터리}

\begin{dhkcode}
Path.home()
\end{dhkcode}

Mac에서는 예를 들어:

\begin{dhkzsh}
/Users/user
\end{dhkzsh}

같은 결과가 나온다.

\subsubsection{경로 결합}

가장 자주 쓴다.

\begin{dhkcode}
project = Path.cwd()
data_dir = project / "data"
train_file = data_dir / "train.csv"
\end{dhkcode}

또는 한 번에:

\begin{dhkcode}
train_file = Path.cwd() / "data" / "train.csv"
\end{dhkcode}

개념:

\begin{dhksmry}[centered]
CWD
\da
my_ml_project
\da
data
\da
train.csv
\end{dhksmry}

\subsubsection{파일/디렉터리 존재 여부}

\begin{dhkcode}
path = Path("data/train.csv")

print(path.exists())
\end{dhkcode}

존재하면 \zsh{True}, 없으면 \zsh{False}

ML 코드에서 굉장히 유용하다.

\begin{dhkcode}
if not path.exist():
    print("데이터 파일이 없습니다.")
\end{dhkcode}

\subsubsection{파일인지 확인}

\begin{dhkcode}
path.is_file()
\end{dhkcode}

예:

\begin{dhkcode}
Path("data/train.csv").is_file()
\end{dhkcode}

실제 파일이면 \zsh{True}, 아니면 \zsh{False}

\subsubsection{디렉터리인지 확인}

\begin{dhkcode}
Path("data").is_dir()
\end{dhkcode}

실제 디렉터리면 \zsh{True}, 아니면 \zsh{False}

따라서:

\begin{dhksmry}[centered]
exists()
\ra 존재하는가?

is_file()
\ra 파일인가?

is_dir()
\ra 디렉터리인가?
\end{dhksmry}

로 구분하면 된다.

\subsubsection{파일 이름 가져오기}

\begin{dhkcode}
path = Path("/Users/user/Desktop/project/data/train.csv")

print(path.name)
\end{dhkcode}

결과:

\begin{dhkzsh}
train.csv
\end{dhkzsh}

\subsubsection{확장자 가져오기}

\begin{dhkcode}
print(path.suffix)
\end{dhkcode}

결과:

\begin{dhkzsh}
.csv
\end{dhkzsh}

ML/DL에서는:

\begin{dhkzsh}
.csv
.json
.pt
.pth
.ckpt
.safetensors
\end{dhkzsh}

같은 확장자를 다룰 일이 많다.

\subsubsection{확장자를 제외한 파일 이름}

\begin{dhkcode}
print(path.stem)
\end{dhkcode}

결과:

\begin{dhkzsh}
train
\end{dhkzsh}

즉:


\begin{dhkdiagram}[in=10]
node stem (1,1) :: train.csv
node sd (3,1)  :: stem = train
node fd (2,3)  :: suffix = .csv

edge stem -> sd
edge stem -> fd :: VH
\end{dhkdiagram}

그리고:

\begin{dhkzsh}
name = train.csv
\end{dhkzsh}

이다.

\subsubsection{부모 디렉터리}

\begin{dhkcode}
print(path.parent())
\end{dhkcode}

경로가:

\begin{dhkzsh}
/Users/user/project/data/train.csv
\end{dhkzsh}

라면:

\begin{dhkzsh}
/Users/user/project/data
\end{dhkzsh}

가 나온다.

다시:

\begin{dhkcode}
path.parent.parent
\end{dhkcode}

하면:

\begin{dhkzsh}
/Users/user/project
\end{dhkzsh}

이다.

아까 \zsh{__file__}과 결합했던 이유가 이것이다.

\subsubsection{확장자 변경}

\begin{dhkcode}
path = Path("data/train.csv")
new_path = path.with_suffix(".json")
print(new_path)
\end{dhkcode}

결과:

\begin{dhkzsh}
data/train.json
\end{dhkzsh}

주의할 점은 이 코드가 \textbf{실제 파일의 확장자를 변경한 것은 아니다}.

새로운 경로 객체를 만든 것이다.

이 구분은 중요하다.

\begin{dhkcode}
path.with_suffix(".json")
\end{dhkcode}

은:
\begin{dhksmry}[centered]
.json으로 끝나는 새로운 Path를 만들어라.
\end{dhksmry}

이지, \textbf{SSD의 파일 이름을 실제로 바꿔라가 아니다}.

\subsubsection{파일 이름 자체 변경}

경로 표현만 바꾸려면:

\begin{dhkcode}
path = Path("data/train.csv")
new_path = path.with_name("train_v2.csv")
\end{dhkcode}

결과:

\begin{dhkzsh}
data/train_v2.csv
\end{dhkzsh}

이것 역시 \textbf{경로 객체 생성}이다.

실제 파일 변경과는 다르다.

\subsubsection{실제 디렉터리 만들기}

이제부터는 실제 파일 시스템에 영향을 준다.

\begin{dhkcode}
output_dir = Path("output")
output_dir.mkdir()
\end{dhkcode}

실제로 \zsh{output/}디렉터리가 생성된다.

이미 있으면 오류가 날 수 있으므로 보통:

\begin{dhkcode}
output_dir.mkdir(exist_ok=True)
\end{dhkcode}

라고 한다.

\subsubsection{중첩 디렉터리 생성}

\begin{dhkcode}
checkpoint_dir = Path("outputs/models/checkpoints")

checkpoint_dir.mkdir(
    parents = True,
    exist_ok = True
)
\end{dhkcode}

그러면 필요한 부모 디렉터리까지 만든다.

\begin{dhkdir}
outputs/
>models/
>checkpoints/
\end{dhkdir}

ML 코드에서 정말 자주 보는 패턴이다.

\subsubsection{간단한 텍스트 파일 쓰기}

\zsh{Path}자체로도 가능하다.

\begin{dhkcode}
path = Path("result.txt")

path.write_text("accuracy: 0.95")
\end{dhkcode}

실제로 \zsh{result.txt}가 생성된다.

읽으려면:

\begin{dhkcode}
text = path.read_text()
print(text)
\end{dhkcode}

결과:

\begin{dhkzsh}
accuracy: 0.95
\end{dhkzsh}

\subsubsection{파일 삭제}

\begin{dhkcode}
path = Path("result.txt")
path.unlink()
\end{dhkcode}

실제 파일을 삭제한다.

따라서 \zsh{unlink()}는 앞에서 본:

\begin{dhkcode}
with_suffix(x)
with_name()
\end{dhkcode}

보다 훨씬 주의해야 한다.

\zsh{unlink()}는 실제 파일 시스템을 변경하기 때문이다.

\subsubsection{빈 디렉터리 삭제}

\begin{dhkcode}
directory = Path("outputs")
directory.rmdir()
\end{dhkcode}

단, 기본적으로 비어있는 디렉터리여야 한다.

파일이 여러 개 들어 있는 폴더를 통째로 제거하는 것은 별도의 문제로 다루자.

\subsubsection{디렉터리 안을 살펴보기}

프로젝트가:

\begin{dhkdir}
data/
>train.csv
>labels.json
>README.txt
\end{dhkdir}

라고 하자.

\begin{dhkcode}
data_dir = Path("data")

for path in data_dir.iterdir():
    print(path)
\end{dhkcode}

대략:

\begin{dhkzsh}
data/train.csv
data/test.csv
data/labels.json
data/README.txt
\end{dhkzsh}

가 나온다.

즉 \zsh{iterdir()}는:

\begin{dhksmry}[centered]
이 디렉터리 바로 아래 뭐가 있는지 순회
\end{dhksmry}

한다.

\subsubsection{특정 파일만 찾기 - \zsh{glob}}

csv만 찾고 싶다.

\begin{dhkcode}
data_dir = Path("data")
for path in data_dir.glob("*.csv"):
    print(path)
\end{dhkcode}

결과:

\begin{dhkzsh}
data/train.csv
datat/test.csv
\end{dhkzsh}

여기서:

\begin{dhkzsh}
*.csv
\end{dhkzsh}

는:

\begin{dhksmry}[centered]
이름은 무엇이든 상관없고 .csv로 끝나는 것
\end{dhksmry}

이라는 패턴이다.

*를 wildcard라고 부른다.

\subsubsection{하위 디렉터리까지 전부 찾기 - \zsh{rglob}}

구조가:

\begin{dhkdir}
data/
>train.csv
>raw/
>>raw_data.csv
>processed/
>>clean_data.csv
\end{dhkdir}

라고 하자.

\begin{dhkcode}
for path in Path("data").rglob("*.csv"):
    print(path)  
\end{dhkcode}

그러면 하위 디렉터리까지 재귀적으로 검색해서:

\begin{dhkzsh}
data/train.csv
data/raw_data.csv
data/processed/clean_data.csv
\end{dhkzsh}

를 찾을 수 있다.

ML 데이터셋 관리에서 꽤 유용하다.

\subsubsection{이름 변경 - 실제 파일 시스템 조작}

아까:

\begin{dhkcode}
with_name()
\end{dhkcode}

은 새로운 경로만 만들었다.

실제 파일 이름을 변경하려면 예를 들어:

\begin{dhkcode}
old = Path("data/train.csv")
new = Path("data/train_v2.csv")

old.rename(new)
\end{dhkcode}

라고 할 수 있다.

실제로:

\begin{dhksmry}[centered]
train.csv \ra train_v2.csv
\end{dhksmry}

로 변경된다.

이 차이를 기억하자.

\begin{dhksmry}[centered]
with_name()
with_suffix()
\da
Path 표현 생성

rename()
unlink()
mkdir()
\da
실제 파일 시스템 변경
\end{dhksmry}

\subsubsection{\zsh{Path}에서 아주 중요한 개념적 구분}

이걸 객체 관점에서 보면 더 명확하다.

\begin{dhkcode}
p = Path("data/train.csv")
\end{dhkcode}

여기서 \zsh{p}는 \textbf{파일 그 자체가 아니다}.

\zsh{p}는:

\begin{dhksmry}[centered]
파일 시스템상의 어떤 위치를 표현하는 Python 객체
\end{dhksmry}

이다.

그래서:

\begin{dhkcode}
p.exists()
\end{dhkcode}

라고 물어보는 것이다.

즉:

\begin{dhkbranch}[in=8]
node main :: Path 객체
arrow main :: 이 위치에
node main :: data/train.csv
arrow main :: 실제 파일 시스템을 확인
node main :: 파일이 있는가?
\end{dhkbranch}

이다.

이 관점은 Python의 \textbf{name $\rightarrow$ object} 사고방식과도 잘 맞는다.

\begin{dhkcode}
p = Path("data/train.csv")
\end{dhkcode}

에서 \zsh{p}라는 이름이 \zsh{Path}객체를 가리키고 있는 것이다.

그리고 그 \zsh{Path}객체가 파일 시스템상의 위치를 표현한다.

\subsubsection{실전 ML 코드로 합쳐보자}

예를 들어:

\begin{dhkcode}
from pathlib import Path
import pandas as pd

project_root = Path(__file__).resolve().parent.parent

data_dir = project_root / "data"
model_dir = project_root / "models"

train_path = data_dir / "train.csv"

if not train_path.exists():
    raise FileNotFoundError(f'데이터가 없습니다: {train_path}')

mdoel_dir.mkdir(
    parents = True,
    exist_ok = True
)

df = pd.read_csv(train_path)
\end{dhkcode}

이제 이 코드가 거의 전부 읽혀야 한다.

구조는:

\begin{dhkbranch}[in=8]
node main :: 현재 Python 파일
arrow main
node main :: __file__
arrow main
node main :: .parent.parent
arrow main
node main :: project root

branchline main -> data :: data //* train.csv
stop data

handoff main -> models :: models
stop models
\end{dhkbranch}

그리고:

\begin{dhksmry}[centered]
trian_path.exists()
\da
데이터 존재 확인

model_dir.mkdir()
\da
모델 저장 디렉터리 준비

pd.read_csv(train_path)
\da
실제 데이터 읽기
\end{dhksmry}

가 된다.

\subsubsection{마지막으로 세 개를 확실히 분리하자}

프로젝트:

\begin{dhkdir}
my_ml_project
>data/
>>train.csv
>src/
>>train.py
\end{dhkdir}

프로젝트 root에서:

\begin{dhkzsh}
python src/train.py
\end{dhkzsh}

를 실행했다.

\begin{dhknote}
\begin{dhksmry}
Path.cwd()
\ra /Users/user/Desktop/my_ml_project
의미: \textbf{내 작업 기준점은 어디인가?}

Path(__file__).resolve()
\ra /Users/user/Desktop/my_ml_project/src/train.py
의미: \textbf{현재 이 Python 파일은 어디 있는가?}

sys.path
\ra [
    ...,
    ".../site-packages",
    ...
    ]
의미: \textbf{Python이 import할 모듈/패키지를 어디에서 찾는가?}
\end{dhksmry}
\end{dhknote}

이 세 가지는 \textbf{모두 경로를 다루지만 서로 다른 질문에 대한 답}이다.

\subsubsection{여기서 다음 단계는 \zsh{sys.argv}이다}

이제 드디어 \textbf{커맨드라인}으로 들어갈 기반이 거의 만들어졌다.

다음에는 이 한 줄을 완전히 해부하면 된다.

\begin{dhkzsh}
python src/train.py --data data/train.csv --epochs 10
\end{dhkzsh}

지금까지 배운 걸 기준으로 이미 앞부분은 이해할 수 있다.

\begin{dhktree}
python src/train.py --data data/train.csv --epoch 10
>python :: 파이썬 인터프리터
>src/train.py :: 실행할 파일
>data.train.csv :: argument
>10 :: argument 

\end{dhktree}

터미널에:

\begin{dhkzsh}
python src/train.py --data data/train.csv --epochs 10
\end{dhkzsh}

을 입력한다.

이걸 구조적으로 보면:

\begin{dhkbranch}[in=8]
node main :: 사용자
arrow main :: 키보드 입력
node main :: 터미널
arrow main
node main :: Shell
arrow main :: 명령 해석
node main :: Python interpreter 실행
arrow main :: train.py 실행
node main :: Python 프로그램
arrow main :: sys.argv
node main :: 커맨드 라인에서 전달된 값 사용
\end{dhkbranch}

여기서 새로운 개념들이 등장한다.

\begin{dhksmry}[centered]
Terminal
Shell
Command
Argument
Option
Python interpreter
sys.argv
\end{dhksmry}

하나씩 분리하자.

\subsubsection{터미널(Terminal)은 무엇인가?}

Mac에서 Terminal, iTerm, VS Code의 Terminal 같은 화면을 본다고 하자.

예를 들어:

\begin{dhkzsh}
(transformer) user@MacBook my_ml_project %
\end{dhkzsh}

그리고 여기에:

\begin{dhkzsh}
python src/train.py
\end{dhkzsh}

라고 입력한다.

흔히 이걸 보고:

\begin{dhksmry}[centered]
터미널이 Python을 실행한다.
\end{dhksmry}

라고 말하지만 조금 더 정확하게 보면 \textbf{터미널과 Shell은 다르다}.

터미널은 사용자와 Shell이 상호작용할 수 있는 \textbf{입출력 인터페이스}에 가깝다.

쉽게 말하면:
\begin{dhkzsh}
Terminal
---------------------------
python train.py 
hello
---------------------------
\end{dhkzsh}

처럼 명령을 입력하고 결과를 볼 수 있는 창을 제공한다.

\subsubsection{실제로 명령을 해석하는 것은 Shell}

Mac에서는 흔히 zsh를 사용한다.

Linux에서는 zsh, bash 등이 있고,

Windows에서는 PowerShell, cmd 등을 만날 수 있다.

Shell의 역할은 대략:

\begin{dhksmry}[centered]
사용자가 입력한 명령을 해석해서 적절한 프로그램을 실행하는 것
\end{dhksmry}

이다.

그래서:

\begin{dhkzsh}
python src/train.py
\end{dhkzsh}

라고 입력하면 Shell이 이것을 해석한다.

\subsubsection{공백이 상당히 중요하다}

다음 명령을 보자.

\begin{dhkzsh}
python src/train.py --data data/train.csv --epochs 10
\end{dhkzsh}

Shell입장에서는 대략 공백을 기준으로 여러 개의 토큰으로 분리된다.

\begin{dhkzsh}
python
src/train.py
--data
data/train.csv
--epochs
10
\end{dhkzsh}

즉 개념적으로:

\begin{dhkzsh}
[python]
[src/train.py]
[--data]
[data/train.csv]
[--epochs]
[10]
\end{dhkzsh}

처럼 된다.

물론 따옴표, escape 등의 규칙 때문에 실제 Shell parsing은 이것보다 더 복잡하지만 지금은 이정도면 충분하다.

\subsubsection{첫 번째 \zsh{python}은 무엇인가?}

\begin{dhkdir}
python src/train.py
> 이것은 실행하려는 command다.
\end{dhkdir}

Shell입장에서는:

\begin{dhksmry}[centered]
python이라는 프로그램을 실행하라.
\end{dhksmry}

라는 의미다.

여기서 Python interpreter가 등장한다.

\subsubsection{Python interpreter란?}

우리가 작성하는:

\begin{dhkcode}
print("hello")
\end{dhkcode}

는 그 자체로 CPU가 바로 실행하는 기계어가 아니다.

Python 코드를 읽고 실행하는 프로그램이 필요하다.

그 프로그램이 \textbf{Python interpreter}다.

따라서:

\begin{dhkzsh}
python src/train.py
\end{dhkzsh}

는 대략:

\begin{dhksmry}[centered]
Shell
"python이라는 프로그램을 실행해."

\da

Python interpreter
"좋아. 그런데 뭘 실행하지?"

\da

src/train.py
\end{dhksmry}

라는 구조다.

\subsubsection{Conda와도 바로 연결된다}

예를 들어:

\begin{dhkzsh}
conda activate transformer
\end{dhkzsh}

를 했다.

프롬프트가:

\begin{dhkzsh}
(transformer) %
\end{dhkzsh}

처럼 바뀌었다.

그리고:

\begin{dhkzsh}
python
\end{dhkzsh}

을 실행한다.

여기서 중요한 질문은:

\begin{dhksmry}[centered]
수많은 Python 중 어떤 Python이 실행되는가?
\end{dhksmry}


이다.

Conda 환경이 활성화되어 있다면 보통 해당 환경의 Python interpreter가 선택된다.

Windows transformer 환경을 예로 들면:

\begin{dhkzsh}
C:\Users\user\anaconda3\envs\transformer\python.exe
\end{dhkzsh}

같은 실행 파일이다.

Python 안에서도 확인할 수 있다.

\begin{dhkcode}
import sys

print(sys.executable)
\end{dhkcode}

즉:

\begin{dhkbranch}[in=8]
node main :: Shell
arrow main :: python
node main :: Conda transformer 환경의 python.exe
arrow main
node main :: Python interpreter
\end{dhkbranch}

라고 연결된다.

\subsubsection{여기서 \zsh{sys.executable}의 정체가 나온다}

이제 \zsh{sys}가 등장한다.

\begin{dhkcode}
import sys

print(sys.executable)
\end{dhkcode}

\zsh{sys.executable}은 대략:

\begin{dhknote}
\begin{dhksmry}[centered]
현재 이 Python 프로그램을 실행하고 있는

Python interpreter 실행 파일은 어디 있는가?
\end{dhksmry}
\end{dhknote}

를 알려준다.

예를 들어:

\begin{dhkzsh}
C:\Users\user\anaconda3\envs\transformer\python.exe
\end{dhkzsh}

그러니까 이것도 단순한 파일 경로 이야기가 아니다.

\textbf{현재 실행 중인 Python 환경을 확인하는 것}이다.

그래서 환경 문제를 디버깅할 때 굉장히 유용하다.

예를 들어:

\begin{dhksmry}[centered]
분명 transforemrs를 설치했는데 import transformers가 안되는데?
\end{dhksmry}

이럴 때:

\begin{dhkcode}
import sys

print(sys.executable)
\end{dhkcode}

을 보면 엉뚱한 Python을 실행하고 있었다는 걸 발견하는 경우가 있다.

\subsubsection{이제 두 번재 토큰 \zsh{src/train.py}}

전체 명령:

\begin{dhkzsh}
python src/train.py --data data/train.csv --epochs 10
\end{dhkzsh}

여기에서:

\begin{dhkdir}
python
>src/train.py
\end{dhkdir}

이다.

Shell은 Python 프로그램을 실행하면서 나머지 문자열들을 Python 쪽에 전달한다.

Python interpreter 입장에서 첫 번째 중요한 값은:

\begin{dhkzsh}
src/train.py
\end{dhkzsh}

이다.

즉:

\begin{dhksmry}
이 Python 스크립트를 실행해라.
\end{dhksmry}

\subsubsection{여기서 앞에서 배운 CWD가 바로 등장한다}

현재 터미널 위치가:

\begin{dhkzsh}
my_ml_project/
\end{dhkzsh}

라고 하자.

구조:

\begin{dhkdir}
my_ml_project/
>data/
>>train.csv
>src/
>>train.py
\end{dhkdir}

그리고:

\begin{dhkzsh}
python src/train.py
\end{dhkzsh}

를 실행한다.

\zsh{src/train.py}는 \textbf{상대 경로}다.

무엇을 기준으로?

앞에서 배운 CWD다.

\begin{dhksmry}[centered]
CWD
\da
my_ml_project
+
src/train.py
\da
my_ml_project/sr/train.py
\end{dhksmry}

앞에서 공부한 파일 경로와 커맨드 라인이 연결되었다.

\subsubsection{그런데 뒤에 이상한 것들이 더 있다}

\begin{dhkzsh}
python src/train.py --data data/train.csv --epochs 10
\end{dhkzsh}

Python 파일 뒤에:

\begin{dhkzsh}
--data
data.train.csv
--epochs
10
\end{dhkzsh}

이 있다.

이것들이 \zsh{커맨드라인 인수(command-line arguments)}다.

프로그램을 실행하면서 프로그램에 추가 정보를 전달하는 것이다.

\subsubsection{왜 굳이 커맨드라인으로 정보를 전달할까?}

예를 들어 \zsh{train.py}에 이렇게 작성할 수도 있다.

\begin{dhkcode}
data_path = "data/train.csv"
epochs = 10
\end{dhkcode}

그런데 다음번에는 다른 데이터를 학습하거나 epochs를 변경하고 싶다.

그러면:

\begin{dhkcode}
data_path = "data/test.csv"
epochs = 20
\end{dhkcode}

으로 수정해야 한다.

불편하다.

그래서 프로그램 자체와 실행 설정을 분리한다.

\begin{dhkzsh}
python src/train.py --data data/train.csv --epochs 10
\end{dhkzsh}

다음에는:

\begin{dhkzsh}
python src/train.py --data data/new_data.csv --epochs 30
\end{dhkzsh}

라고 실행할 수 있다.

\zsh{train.py} 코드는 수정하지 않았다.

\subsubsection{이게 커맨드라인 인수의 핵심 목적이다}

\begin{dhksmry}[centered]
프로그램 코드
train.py

+

실행할 때 필요한 설정
data/train.csv
epochs = 10
\end{dhksmry}

를 분리한다.

즉:

\begin{dhkdir}[in=8]
train.py
>실행 A
>>data=train.csv
>>epochs=10
>실행 B
>>data=train2.csv
>>epochs=20
>실행 C
>>data=train3.csv
>>epochs=50
\end{dhkdir}

가 가능해진다.

ML/DL에서 CLI를 많이 사용하는 이유 중 하나다.

\subsubsection{이제 Argument와 Option을 구분하자}

간단한 명령:

\begin{dhkzsh}
python train.py data/train.csv
\end{dhkzsh}

여기서:

\begin{dhkzsh}
data/train.csv
\end{dhkzsh}

는 argument다.

프로그램에 전달하는 값이다.

그런데 값이 많아지면 문제가 생긴다.

\begin{dhkzsh}
python train.py data/train.csv 10 32 model.pt
\end{dhkzsh}

이게 뭘 의미하는지 알아보기 어렵다.

\begin{dhksmry}[centered]
data/train.csv \ra 데이터?

10 \ra epochs?

32 \ra batch size?

model.pt \ra outputs?
\end{dhksmry}

그래서 이름을 붙여준다.

\begin{dhkzsh}
python train.py \
    --data data/train.csv \
    --epochs 10 \
    --batch-size 32 \
    --output model.pt
\end{dhkzsh}

훨씬 명확하다.

\subsubsection{\zsh{--data, --epochs}가 option이다}

보통:

\begin{dhkzsh}
--data
--epochs
--batch-size
--output
\end{dhkzsh}

같은 것을 option이라고 부른다.

그리고:

\begin{dhkzsh}
--data data/train.csv
\end{dhkzsh}

를 보면:

\begin{dhkdir}
--data
> 어떤 설정인지 나타내는 option
\end{dhkdir}
\begin{dhkdir}
data/train.csv
> 전달할 value
\end{dhkdir}


라고 볼 수 있다.

다만 CLI 관련 용어는 라이브러리/문서에 따라 

argument, option, parameter 등을 조금 다르게 쓰기도 한다. 

중요한 것은 구조다.

\subsubsection{위치로 의미를 판단하는 positional argument도 있다}

예를 들어:

\begin{dhkzsh}
python copy.py source.txt destination.txt
\end{dhkzsh}

이라고 하자.

프로그램이:

\begin{dhksmry}[centered]
첫 번째 값 = source

두 번째 값 = destination
\end{dhksmry}

이라고 약속했다면:

\begin{dhkzsh}
source.txt
destination.txt
\end{dhkzsh}

는 위치 자체가 의미를 가진다.

이것을 \textbf{positional argument라고 한다}.

예를 들어:

\begin{dhkzsh}
git add README.md
\end{dhkzsh}

같은 명령에서도 비슷한 사고방식을 볼 수 있다.

\subsubsection{Option은 이름을 붙여 전달한다}

반면:

\begin{dhkzsh}
python train.py --epochs 10
\end{dhkzsh}

에서는:

\begin{dhkzsh}
--epochs
\end{dhkzsh}

라는 이름을 통해 10의 의미를 알 수 있다.

그래서 개념적으로:

\begin{dhksmry}[centered]
Positional - 위치로 의미 결정

python train.py data.csv 10
\end{dhksmry}

위의 내용 대비:

\begin{dhksmry}[centered]
Option - 이름으로 의미 명시

python train.py --data data.csv --epochs 10
\end{dhksmry}

라고 생각하면 된다.

\subsubsection{이제 드디어 \zsh{sys.argv}}

Python은 자신에게 전달된 커맨드라인 인수에 접근할 방법이 필요하다.

그것이:

\begin{dhkcode}
import sys

print(sys.argv)
\end{dhkcode}

이다.

\zsh{argv}는 전통적으로:

\begin{dhksmry}[centered]
argument vector
\end{dhksmry}

에서 온 이름이다.

여기서 vector를 NumPy 벡터처럼 수학적으로 생각할 필요는 없다.

\textbf{전달된 argument들의 순서 있는 목록}정도로 이해하면 충분하다.

\subsubsection{직접 실험해보자}

\zsh{train.py}:

\begin{dhkcode}
import sys

print(sys.argv)
\end{dhkcode}

터미널:

\begin{dhkzsh}
python train.py
\end{dhkzsh}

결과는 대략:

\begin{dhkzsh}
['train.py']
\end{dhkzsh}

여기서 중요한 사실이 나온다.

\zsh{sys.argv[0]}에는 보통 실행 대상 스크립트 이름/경로가 들어간다.

\subsubsection{인수를 하나 전달해보자}

\begin{dhkzsh}
python train.py hello
\end{dhkzsh}

그러면:

\begin{dhkzsh}
['train.py', 'hello']
\end{dhkzsh}

가 된다.

즉:

\begin{dhkdir}
sys.argv
>[0] = 'train.py'
>[1] = 'hello'
\end{dhkdir}

이다.

\subsubsection{여러 개 전달}

\begin{dhkzsh}
python train.py hello world 10
\end{dhkzsh}

그러면:

\begin{dhkzsh}
[
    'train.py',
    'hello',
    'world',
    '10'
]
\end{dhkzsh}

가 된다.

인덱스로 접근이 가능하다.

\begin{dhkcode}
import sys

print(sys.argv[0])
print(sys.argv[1])
print(sys.argv[2])
print(sys.argv[3])
\end{dhkcode}

결과

\begin{dhkzsh}
train.py
hello
world
10
\end{dhkzsh}

\subsubsection{아주 중요한 점: 전부 문자열이다}

\begin{dhkzsh}
python train.py 10
\end{dhkzsh}

이라고 했다고 하자.

사람은 10을 보고 숫자라고 생각한다.

하지만:

\begin{dhkcode}
import sys

print(sys.argv[1])
print(type(sys.argv[1]))
\end{dhkcode}

결과:

\begin{dhkzsh}
10
<class 'str'>
\end{dhkzsh}

이다.

즉:

\begin{dhksmry}[centered]
터미널에서 전달
10
\da
sys.argv
\da
"10"
\end{dhksmry}

이다.

\textbf{커맨드라인 인수는 문자열 형태로 들어오므로 필요하면 변환해야 한다}.

\subsubsection{epochs라면 숫자로 변환해야 한다}

\begin{dhkzsh}
python train.py 10
\end{dhkzsh}

\zsh{train.py}:

\begin{dhkcode}
import sys

epochs = int(sys.argv[1])

print(epochs)
pirnt(type(epochs))
\end{dhkcode}

결과:

\begin{dhkzsh}
10
<class 'int'>
\end{dhkzsh}

즉 \zsh{input()}과 비슷한 면이 있다.

\begin{dhkzsh}
x = input()
\end{dhkzsh}

도 기본적으로 문자열을 받는다.

커맨드라인 인수도:

\begin{dhkzsh}
sys.argv[1]
\end{dhkzsh}

에서 문자열을 받는다.

\subsubsection{파일 경로와 결합해보자}

터미널:

\begin{dhkzsh}
python train.py data/train.csv
\end{dhkzsh}

Python:

\begin{dhkcode}
import sys
from pathlib import Path

data_path = Path(sys.argv[1])

print(data_path)
print(data_path.resolve())
\end{dhkcode}

여기서 흐름이 중요해진다.

\begin{dhkbranch}[in=8]
node main :: Shell
arrow main :: "data/train.csv"

node main :: Python
arrow main :: sys.argv[1]

node main :: "data/train.csv"
arrow main :: Path(...)

node main :: Path("data/train.csv")
arrow main :: CWD 기준

node main :: 실제 파일 시스템상의 위치
\end{dhkbranch}

지금가지 공부한 두 파트가 완전히 연결되었다.

\subsubsection{Pandas까지 붙이면}

\begin{dhkcode}
import sys
import pandas as pd
from pathlib import Path

data_path = Path(sys.argv[1])

df = pd.read_csv(data_path)
print(df.head())
\end{dhkcode}

실행:

\begin{dhkzsh}
python train.py data/train.csv
\end{dhkzsh}

그러면:

\begin{dhkbranch}[in=8]
node main :: Terminal
arrow main :: data/train.csv
node main :: sys.argv
arrow main
node main :: 파일 시스템
arrow main
node main :: train.csv
arrow main
node main :: Pandas
arrow main
node main :: DataFrame
\end{dhkbranch}

이 된다.

이제 CLI가 ML 코드와 연결되기 시작한다.

\subsubsection{그런데 \zsh{sys.argv}만 직접 쓰면 금방 불편해진다}

예를 들어:

\begin{dhkzsh}
python train.py data/train.csv 10 32 model.pt
\end{dhkzsh}

코드는:

\begin{dhkcode}
data = sys.argv[1]
epochs = int(sys.argv[2])
batch_size = int(sys.argv[3])
outuput = sys.argv[4]
\end{dhkcode}

가 된다.

문제가 많다.

순서를 기억해야 하고,

\begin{dhkzsh}
python train.py data/train.csv 32 10 model.pt
\end{dhkzsh}

처럼 순서를 잘못 넣어도 사람이 알아채기 어렵다.

인수가 빠지면 \zsh{IndexError}가 날 수도 있고 도움말도 없다.

\subsubsection{그래서 \zsh{argparse}가 등장한다}

Python 표준 라이브러리에:

\begin{dhkcode}
import argparse
\end{dhkcode}

가 있다.

이걸 사용하면:

\begin{dhkzsh}
python train.py \
    --data data/train.csv \
    --epochs 10 \
    --batch-size 32
\end{dhkzsh}

같은 CLI를 훨씬 체계적으로 만들 수 있다.

예를 들어:

\begin{dhkcode}
import argparse

parser = argparse.ArgumentParser()

parser.add_argument("--data")
parser.add_argument("--epochs, type=int)

args = parser.parse_args()

print(args.data)
print(args.epochs)
\end{dhkcode}

실행:

\begin{dhkzsh}
python train.py --data data/train.csv --epochs 10
\end{dhkzsh}

그러면:

\begin{dhkcode}
args.data
\end{dhkcode}

는:

\begin{dhkzsh}
datat/train.csv
\end{dhkzsh}

이고:

\begin{dhkcode}
args.epochs
\end{dhkcode}

는 \zsh{type=int}를 지정했으므로 정수 10 이다.

\subsubsection{\zsh{argparse}가 마법을 쓰는 것은 아니다}

이것도 구조적으로 이해하는 게 중요하다.

밑바닥에는 여전히 커맨드라인에서 전달된 argument가 있다.

개념적으로:

\begin{dhkbranch}[in=4]
node main :: Shell
arrow main
node main :: python train.py --data data/train.csv --epoch 10
arrow main
node main :: command-line arguments
arrow main
node main :: Python 프로세스

branchline main -> argv :: sys.argv //* raw 문자열 목록
stop argv

handoff main -> argparse :: argparse //* 해석/검증/변환
stop argparse
\end{dhkbranch}

\zsh{sys.argv}는 비교적 \textbf{날것 그대로 접근}하는 방법이고,

\zsh{argparse}는 그것을 \textbf{구조적으로 해석해주는 도구}라고 보면 된다.

\subsubsection{Option을 조금 더 제대로 만들어보자}

\begin{dhkcode}
import argparse
from pathlib import Path

parser = argparse.Argumentparser()

parser.add_argument(
    "--data",
    type=Path,
    required=True
)

parser.add_argument(
    "--epochs",
    type=int,
    default=10
)

args = parser.parse_args()
\end{dhkcode}

이제:

\begin{dhkzsh}
python train.py --data data/train.csv
\end{dhkzsh}

라고만 실행해도:

\begin{dhkcode}
args.data
\end{dhkcode}

는 \zsh{Path}객체가 되고:

\begin{dhkcode}
args.epochs
\end{dhkcode}

는 기본값 10이 된다.

\subsubsection{\zsh{required=True}}

\begin{dhkcode}
parser.add_argument(
    "--data",
    required=True
)
\end{dhkcode}

라고 하면:

\begin{dhksmry}[centered]
--data는 반드시 전달해야 한다.
\end{dhksmry}

라는 의미다.

따라서:

\begin{dhkzsh}
python train.py
\end{dhkzsh}

라고 실행하면 \zsh{argparse}가 오류를 알려준다.

직접:

\begin{dhkcode}
if len(sys.argv) < 2:
...
\end{dhkcode}

같은 검사를 만들 필요가 줄어든다.

\subsubsection{\zsh{default}}

\begin{dhkcode}
parser.add_argument(
    "--epochs",
    type=int,
    default=10
)
\end{dhkcode}

는:

\begin{dhksmry}[centered]
사용자가 epochs를 안 주면 10을 사용해라.
\end{dhksmry}

라는 뜻이다.

따라서:

\begin{dhkzsh}
python train.py --data data/train.csv
\end{dhkzsh}

라면:

\begin{dhkcode}
args.epochs == 10
\end{dhkcode}

이고

\begin{dhkzsh}
python train.py --data data/train.csv --epochs 50
\end{dhkzsh}

라면:

\begin{dhkcode}
args.epochs == 50
\end{dhkcode}

이다.

\subsubsection{도움말도 자동으로 만들 수 있다}

\begin{dhkcode}
parser.add_argument(
    "--epochs",
    type=int,
    default=10,
    help="Number of training epochs"
)
\end{dhkcode}

그리고:

\begin{dhkzsh}
python train.py --help
\end{dhkzsh}

라고 실행하면 CLI 사용법을 보여준다.

그래서 제대로 만들어진 Python 프로그램에서:

\begin{dhkzsh}
-h
--help
\end{dhkzsh}

를 자주 볼 수 있는 것이다.

\subsubsection{이제 실제 ML 형태로 만들어보자}

\begin{dhkcode}
import argparse
from pathlib import Path

parser = argparse.ArgumentParser(
    description="Train a classification model"
)

parser.add_argument(
    "--data",
    type=Path,
    required=True,
    help="Path to training data"
)

parser.add_argument(
    "--epochs",
    type=int,
    default=10
)

...
\end{dhkcode}

실행:

\begin{dhkzsh}
python src/train.py \
    --data data/train.csv \
    --epochs 20 \
    --batch-size 64 \
    --output models/run1
\end{dhkzsh}

그러면 프로그램 내부에서는:

\begin{dhksmry}[centered]
args.data
\ra Path("data/train.csv")

args.epochs
\ra 20

args.batch_size
\ra 64

args.output
\ra Path("models/run1")
\end{dhksmry}

처럼 사용할 수 있다.

\subsubsection{이제 이 명령 전체가 읽혀야 한다}

\begin{dhkzsh}
python src/train.py --data data/train.csv --epochs 10
\end{dhkzsh}

앞에서부터 읽자.

\begin{dhksmry}
    python \ra Shell에게 Python interpreter를 실행하라고 요청

    src/train.py \ra Python에게 실행할 스크립트 전다

    --data \ra data라는 option

    data/train.csv \ra 해당 option에 전달되는 값

    --epochs \ra epochs라는 option

    10 \ra 전달되는 값
\end{dhksmry}

그리고 Python 안에서는 원시적으로 \zsh{sys.argv}를 통해 이 정보에 접근할 수 있고,

더 체계적인 프로그램에서는 \zsh{argparse}로 파싱한다.

\subsubsection{지금까지 공부한 파일 경로까지 전부 합치면}

\begin{dhkzsh}
python src/train.py --data data/train.csv --epochs 10
\end{dhkzsh}

을 실행했다고 하자.

\begin{dhkbranch}[in=8]
node main :: Terminal
arrow main
node main :: Shell
arrow main :: command 해석
node main :: python
arrow main :: 어떤 Python?
node main :: Conda 환경의 Python interpreter
arrow main :: sys.executable로 확인 가능
node main :: src/train.py
arrow main :: CWD 기준 상대 경로
\end{dhkbranch}

이 구조가 이번 파트의 핵심이다.

\subsubsection{그리고 \zsh{sys}의 역할도 조금씩 선명해진다}

지금까지 등장한 \zsh{sys} 기능은 두 개다.

\begin{dhksmry}
    sys.executable
    \ra 현재 어떤 Python interpreter가 나를 실행하고 있는가?

    sys.argv
    \ra 이 프로그램을 실행할 때 어떤 command-line argument가 전달되었는가?

    그리고 앞에서 잠깐 본:
    sys.path
    \ra import할 모듈을 어디에서 찾는가?
\end{dhksmry}

따라서 \zsh{sys}라는 모듈의 공통점이 보이기 시작한다.

\begin{dhkdir}
sys
> 현재 실행중인 Python과 실행 환경에 대한 접근
>> executable - interpreter
>> argv - 실행 인수
>> path - import 검색
\end{dhkdir}

\textbf{\zsh{sys}를 파일 경로 라이브러리라고 생각하면 안 되는 이유}도 상당히 명확해진다.

\subsubsection{\zsh{argparse}에 대한 추가 내용}

\begin{dhkcode}
import argparse
\end{dhkcode}

그냥 Python 표준 라이브러리의 \zsh{argparse} 모듈을 가져온 것이다.

우리가 사용하고 싶은 것은 이 안에 있는 \zsh{argparse.ArgumentParser}다.

\begin{dhkcode}
parser = argparse.ArgumentParser()
\end{dhkcode}

여기서 객체 관점으로 보는 게 중요하다.

\begin{dhkbranch}[in=8]
node main :: argparse
line main
node main :: ArgumentParser
arrow main :: ()
node main :: ArgumentParser 객체
arrow main
node main :: parser
\end{dhkbranch}

즉:

\begin{dhkcode}
parser = argparse.ArgumentParser()
\end{dhkcode}

는 \zsh{ArgumentParser}를 호출하여 객체를 만들고 \zsh{parser}라는 이름에 바인딩 한 것이다.

이 객체의 역할은:

\begin{dhksmry}[centered]
커맨드라인 입력을 어떻게 해석할 것인지 관리하는 객체
\end{dhksmry}

라고 생각하면 된다.

\begin{dhkcode}
parser.add_argument("--name")
\end{dhkcode}

이제 \zsh{parser}에게 말하는 것이다.

\begin{dhksmry}[centered]
--name 이라는 입력을 받을 거야.
\end{dhksmry}

그리고:

\begin{dhkcode}
parser.add_argument("--age", type=int)
\end{dhkcode}

는:

\begin{dhksmry}[centered]
--age 라는 입력도 받을 거야. 그리고 이 값은 int로 변환해.
\end{dhksmry}

이다.

그러니까 \zsh{add_argument()}는 이름 그대로다.

\begin{dhksmry}[centered]
"내가 어떤 argument를 받아야 하지?"
\da
add_argument("--name")
\da
add_argument("--age", type=int)
\end{dhksmry}

\textbf{입력 규칙을 \zsh{parser}에 하나씩 등록하는 것}이다.

하지만 아직 실제 값을 읽은 게 아니며 이게 중요하다.

\begin{dhkcode}
parser.add_argument("--name")
parser.add_argument("--age", type=int)
\end{dhkcode}

여기까지 했다고:

\begin{dhkzsh}
Kim
30
\end{dhkzsh}

가 들어온 게 아니다.

그냥:

\begin{dhksmry}[centered]
앞으로 이런 입력을 받을 예정이다.
\end{dhksmry}

라는 규칙만 만들어놓은 것이다.

프로그래밍적으로 표현하면:

\begin{dhksmry}[centered]
프로그램의 입력 설계
\end{dhksmry}

단계다.

이제:

\begin{dhkcode}
args = parser.parse_args()
\end{dhkcode}

가 나온다.

여기서 실제 커맨드라인을 읽고 \zsh{parser}가 등록된 규칙에 따라 해석한다.

우리가:

\begin{dhkzsh}
python hello.py --name Kim --age 30
\end{dhkzsh}

라고 실행했다.

그러면 개념적으로:

\begin{dhkbranch}[in=8]
node main :: --name Kim --age 30
arrow main
node main :: parser.parse_args()

branchline main -> name :: --name 발견
arrow name
node name :: Kim
stop name

handoff main -> age :: --age 발견
arrow age
node age :: "30"
arrow age :: type=int
node age :: 30
stop age
\end{dhkbranch}

그리고 결과 객체를 반환한다.

그 객체에 \zsh{args}라는 이름을 붙인 것이다.

이후:

\begin{dhkcode}
print(args)
\end{dhkcode}

를 해보자.

대략:

\begin{dhkzsh}
Namespace(name='Kim', age=30)
\end{dhkzsh}

가 나온다.

즉 \zsh{parser}가 지저분한 \zsh{command-line}입력을:

\begin{dhkzsh}
--name Kim --age 30
\end{dhkzsh}

에서:

\begin{dhkdir}
Namespace
>name = "Kim"
>age = 35
\end{dhkdir}

라는 Python에서 다루기 쉬운 객체로 정리해 준 것이다.

그래서:

\begin{dhkcode}
args.name
args.age
\end{dhkcode}

로 접근할 수 있다.

이게 \zsh{sys.agrv[1],[2],[3]}를 기억하는 것보다 훨씬 편한 이유다.

\begin{dhknote}
\begin{dhksmry}
1. ArgumentParser()

    Argument + Parser 는 argument를 parse하는 것

    parse : \textbf{구조를 분석해서 의미 있는 형태로 해석한다}정도의 의미

    따라서:

    ArgumentParser = command-line argument를 해석하는 객체

    라고 보면 외우지 않아도 이름 자체가 기능 설명임을 알 수 있다.


2. args = parser.parse_args()      

    이것 또한 마찬가지다

    parse + args 이므로 arguments를 해석하라.


3. argparse는 option의 - 를 attribute의 _ 형태로 만들어준다.

    Command Line            Python

    --batch-size      \ra     args.batch_size   
    
    --learning rate   \ra     args.learning_rate

\end{dhksmry}
\end{dhknote}

\subsubsection{기억할 최소 골격과 자주 나타나는 오류}

당장은 이것만 직접 칠 수 있게 되는 것이 목표다.

\begin{dhkcode}
import argparse

parser = argparse.ArgumentParser()
parser.add_argument("--name")
parser.add_argument("--age", type=int)
args = parser.parse_args()
print(args.name)
\end{dhkcode}

그리고 터미널:

\begin{dhkzsh}
python example.py --name Kim --age 30
\end{dhkzsh}

이 작은 프로그램을 \textbf{보지 않고 한 번 작성해보는 것}을 추천한다.

지금부터 몇 번 직접 타이핑하면서 아래 패턴을 손에 넣으면 된다.

\begin{dhksmry}[centered]
만든다
ArgumentParser()

\da

등록한다
add_argument()

\da

해석한다
parser_args()

\da

사용한다
args.xxx
\end{dhksmry}

\begin{dhksmry}[centered]
\zsh{ \ } 는 단순히 줄바꿈 기호가 아니다.
\end{dhksmry}

zsh/bash 같은 Shell에서 줄 끝의 

\begin{dhkcode}
    \
\end{dhkcode}

는 \textbf{다음 줄까지 현재 명령이 계속된다는 뜻}이다.

좀 더 정확하게 말하면 바로 뒤의 newline을 escape해서

그 줄바꿈을 명령 구분으로 취급하지 않게 만든다.

그리고 이 \zsh{ \ }를 작성하며 나오는 \zsh{ > } 또한 오류가 아니다.

터미널에서:

\begin{dhkzsh}
(transformer) ... % python train.py \
>
\end{dhkzsh}

처럼 \zsh{ > }가 나온다면, 이건 zsh가:

\begin{dhksmry}[centered]
명령이 아직 끝나지 않았으니 다음 줄을 입력하세요.
\end{dhksmry}


라고 기다리고 있는 것이다.

\zsh{ > }를 직접 입력하는 것은 아니다.

Shell이 표시해주는 continuation prompt다.

Python REPL의 >>>와 혼동하면 안된다.

\subsubsection{지금까지 배운 네 가지를 정리하고 추가 내용}

지금까지 배운 각각은:

\begin{dhksmry}[centered]
type
\ra 어떤 Python 타입으로 변환할까?

default
\ra 사용자가 입력하지 않으면 무엇을 사용할까?

required
\ra 반드시 입력해야 하는가?

help
\ra --help에서 어떻게 설명할까?
\end{dhksmry}

이다.

그런데 ML/DL 프로그램을 만들다 보면 이것만으로 부족하다.

특히 자주 만날 것은

\begin{dhkzsh}
action
choices
nargs
짧은 option alias
dest
positional argument
\end{dhkzsh}

정도다.

\subsubsection{\zsh{action} - 값을 받지 않는 스위치}

ML에서 이런 옵션을 만들고 싶다고 하자.

\begin{dhkzsh}
python train.py --use-augmentation
\end{dhkzsh}

뜻은:
\begin{dhksmry}[centered]
augmentation을 사용한다.
\end{dhksmry}

여기서는 굳이 \zsh{--use-augmentation True}라고 할 필요가 없다.

\zsh{--use-augmentation}이 \textbf{존재하는 것 자체를 True}로 만들 수 있다.

\begin{dhkcode}
parser.add_argument(
    "--use-augmentation",
    action="store_true"
)
\end{dhkcode}

이제:

\begin{dhkzsh}
python train.py
\end{dhkzsh}

라면:

\begin{dhkcode}
args.use_augmentation
\end{dhkcode}

은 \zsh{False}다.

반대로:

\begin{dhkzsh}
python train.py --use-augmentation
\end{dhkzsh}

하면:

\begin{dhkcode}
args.use-augmentation
\end{dhkcode}

은 \zsh{True}가 된다.

이런 걸 흔히 flag라고 생각하면 된다.

GPU 관련 설정, logging, augmentation, debug mode 같은 on/off 설정에 잘 어울린다.

\subsubsection{ML에서는 \zsh{choices}를 어디에 쓸까?}

예를 들어:

\begin{dhkcode}
parser.add_argument(
    "--optimizer",
    choices=["adamw", "adam", "sgd"],
    default="adamw"
)

parser.add_argument(
    "--scheduler",
    choices=["linear", "cosine", "constant"],
    default="linear"
)

parser.add_argument(
    "--device",
    choices=["cpu", "cuda", "mps"],
    default="cpu"
)
\end{dhkcode}

같은 식이다.

그러면 CLI 자체가 프로그램이 허용하는 설정을 설명한다.

\subsubsection{\zsh{nargs} - argument 하나에 여러 값 받기}

이번에는 hidden layer 크기를 받고 싶다고 하자.

예를 들어:

\begin{dhksmry}[centered]
512 \ra 256 \ra 128
\end{dhksmry}

사용자가:

\begin{dhkzsh}
python train.py --hidden-sizes 512 256 128
\end{dhkzsh}

라고 입력하게 하고 싶다.

그때:

\begin{dhkcode}
parser.add_argument(
    "--hidden-sizes",
    type=int,
    nargs="+"
)
\end{dhkcode}

라고 할 수 있다.

그러면:

\begin{dhkcode}
args.hidden_sizes
\end{dhkcode}

는:

\begin{dhkzsh}
[512, 256, 128]
\end{dhkzsh}

이 된다.

즉 하나의 option에 여러 값을 전달했다.

\subsubsection{\zsh{nargs="+"}의 +는 무슨 뜻인가?}

여기서는:

\begin{dhksmry}[centered]
하나 이상
\end{dhksmry}

이라는 뜻이다.

그래서:

\begin{dhkzsh}
--hidden-sizes 512
\end{dhkzsh}

도 가능하고:

\begin{dhkzsh}
--hidden-szie 512 256 128
\end{dhkzsh}

도 가능하다.

하지만 값 없이:

\begin{dhkzsh}
--hidden-sizes
\end{dhkzsh}

만 쓰면 안된다.

\subsubsection{짧은 option이름도 같이 만들 수 있다}

예를 들어:

\begin{dhkcode}
parser.add_argument(
    "-e",
    "--epochs",
    type=int,
    default=10
)
\end{dhkcode}

그러면 둘 다 가능하다.

\begin{dhkzsh}
python train.py --epochs 20
python train.py -e 20
\end{dhkzsh}

둘 다:

\begin{dhkcode}
args.epochs == 20
\end{dhkcode}

이다.

짧은 옵션은 자주 사용하는 CLI에서 편리하다.

\subsubsection{\zsh{dest} - 최종 Python attribute 이름을 직접 정하기}

우리가 사용했던:

\begin{dhkcode}
parser.add_argument("--learning-rate")
\end{dhkcode}

는 자동으로:

\begin{dhkcode}
args.learning_rate
\end{dhkcode}

가 된다고 했다.

이 최종 목적지 이름을 \zsh{argparse}에서는 \textbf{destination}, 줄여서 \zsh{dest}라고 한다.

직접 지정할 수도 있다.

\begin{dhkcode}
parser.add_argument(
    "--learning_rate",
    dest="lr,
    type=float
)
\end{dhkcode}

그러면:

\begin{dhkzsh}
python train.py --learning-rate 0.001
\end{dhkzsh}

라고 입력했는데 Python에서는:

\begin{dhkcode}
args.lr
\end{dhkcode}

로 접근한다.

즉:

\begin{dhkdir}
--learning-rate
>dest="lr"
>args.lr
\end{dhkdir}

이다.

\subsubsection{\zsh{dest}는 왜 존재할까?}

CLI에서 사람이 읽기 좋은 이름과 Python 내부에서 쓰기 좋은 이름을 분리할 수 있기 때문이다.

예를 들어 터미널에서는

\begin{dhkzsh}
--learning-rate
\end{dhkzsh}

가 명확하다.

하지만 학습 코드에서는:

\begin{dhkzsh}
lr
\end{dhkzsh}

라는 이름을 사용하고 싶을 수도 있다.

그러면 위와 같은 코드로 변경이 가능하다.

다만 특별한 이유가 없다면 자동 변환되는 기본 동작을 그대로 써도 충분하다.

\subsubsection{Positional argument도 만들 수 있다}

지금까지 전부:

\begin{dhkzsh}
--epochs
--batch-size
--learning-rate
\end{dhkzsh}

처럼 --가 붙었다.

그런데:

\begin{dhkcode}
parser.add_argument("data")
\end{dhkcode}

처럼 --없이 정의하면 positional argument가 된다.

그러면 실행할 때:

\begin{dhkzsh}
python train.py data/train.csv
\end{dhkzsh}

라고 한다.

Python에서는:

\begin{dhkcode}
args.data
\end{dhkcode}

로 사용한다.

즉:

positional argument

\begin{dhktree}
python train.py data/train.csv
>python
>train.py
>data/train.csv :: args.data
\end{dhktree}

optional argument

\begin{dhktree}
python train.py --data data/train.csv
>python
>train.py
>--data data/train.csv :: args.data
\end{dhktree}


이다.

\subsubsection{무엇을 positional로 만들고 무엇을 optional로 만들까?}

정답이 하나로 고정된 것은 아니지만 프로그램 설계의 문제다.

예를 들어 프로그램이:

\begin{dhksmry}[centered]
CSV 하나를 반드시 분석한다.
\end{dhksmry}

라는 도구라면:

\begin{dhkzsh}
python analyze.py data.csv
\end{dhkzsh}

처럼 파일을 positional로 두는 것도 자연스럽다.

반면 ML 학습 프로그램처럼 설정이 많다면:

\begin{dhkzsh}
python train.py \
    --data data/train.csv \
    --epochs 20 \
    --batch-size 64
\end{dhkzsh}

처럼 명시적인 optional으로 두면 읽기 쉽다.

\subsubsection{실제 ML CLI를 한번 읽어보자}

이제 꽤 현실적인 코드를 보자.

\begin{dhkcode}
import argparse
from pathlib import Path

parser = argparse.ArgumentParser(
    description="Train a classification model"
)

parser.add_argument(
    "--data",
    type=Path,
    required=True,
    help="Path to training dataset"
)

parser.add_argument(
    "-e",
    "--epochs",
    type=int,
    default=10
)

parser.add_argument(
    "--optimizer",
    choices=["adam", "adamw", "sgd"],
    default="adamw"
)

parser.add_argument(
    "--hidden-sizes",
    type=int,
    nargs="+"
)

parser.add_argument(
    "--use-augmentation",
    action="store_true"
)
\end{dhkcode}

처음 \zsh{argparse}를 봤을 때는 이 코드가 꽤 낯설었지만

지금은 위에서부터 읽을 수 있다.

\begin{dhkcode}
parser.add_argument(
    "--data",
    type=Path,
    required=True
)
\end{dhkcode}

\begin{dhksmry}[centered]
data 경로는 반드시 받아라. 그리고 Path 객체로 변환해라.
\end{dhksmry}

\begin{dhkcode}
parser.add_argument(
    "-e",
    "--epochs",
    type=int,
    default=10
)
\end{dhkcode}

\begin{dhksmry}[centered]
-e 또는 --epochs 로 받을 수 있고 정수이며 안 주면 10.
\end{dhksmry}

\begin{dhkcode}
parser.add_argument(
    "--optimizer",
    choices=["adam", "adamw", "sgd"],
    default="adamw"
)
\end{dhkcode}

\begin{dhksmry}[centered]
optimizer는 셋 중 하나만 허용하고 기본값은 adamw. 
\end{dhksmry}

\begin{dhkcode}
parser.add_argument(
    "--hidden-sizes",
    type=int,
    nargs="+"
)
\end{dhkcode}

\begin{dhksmry}[centered]
정수를 하나 이상 받아서 여러 값으로 저장.
\end{dhksmry}

\begin{dhkcode}
parser.add_argument(
    "--use-augmentation",
    action="store_true"
)
\end{dhkcode}

\begin{dhksmry}[centered]
option이 등장하면 True가 되는 스위치.
\end{dhksmry}

\subsubsection{\zsh{--help}의 가치도 이제 커진다}

이렇게 잘 정의해놓으면:

\begin{dhkzsh}
python train.py --help
\end{dhkzsh}

만 입력해도 \zsh{argparse}가 프로그램 사용법을 만들어준다.

즉 \zsh{add_argument()}에 우리가 입력한 정보는 단순히 값을 읽는 데만 사용되는 게 아니다.

\begin{dhkbranch}[in=8]
node main :: add_argument()
branchline main -> name :: 입력 이름
stop name
branchline main -> type :: 타입 변환
stop type
branchline main -> default :: 기본값
stop default
branchline main -> need :: 필수 여부
stop need
branchline main -> accept :: 허용 값
stop accept
branchline main -> number :: 값의 갯수
stop number
branchline main -> explain :: 설명
stop explain
arrow main
node main :: ArgumentParser
branchline main -> parse :: 입력 parsing
stop parse
handoff main -> help :: help 생성
stop help
\end{dhkbranch}

그래서 CLI 자체가 하나의 \textbf{입력 인터페이스}가 된다.

\begin{dhktable}
| 바꾸고 싶은 것 | CLI 사용 예시 | argparse에서 필요한 것 |
|:---|:---|:---|
| 학습 횟수 | --epochs 20 | type=int |
| batch 크기 | --batch-size 64 | type=int |
| learning rate | --learning-rate 0.001 | type=float |
| 데이터 위치 | --data train.csv | type=Path |
| 모델 저장 위치 | --output model.pt | type=Path |
| optimizer 종류 | --optimizer adamw | choices=... |
| augmentation 켜기 | --use-augmentation | action="store_true" |
| device 선택 | --device mps | choices=... |
| worker 수 | --num-workers 4 | type=int |
| hidden layer 여러 개 | --hidden-sizes 512 256 | nargs="+", type=int |
\end{dhktable}
